\pagestyle{empty}
\tight
\begin{tblr}{width=\textwidth,colspec={|Q[c,m,1.9cm]|X[l,m]|},hlines,vlines,hspan=minimal,row{1}={bg=headblue},rowsep=1pt,colsep=3pt}
\SetCell{c,m}\hfil\logo\hfil & \textbf{POLITEKNIK NEGERI MALANG}\par\textbf{JURUSAN TEKNOLOGI INFORMASI}\par\textbf{PROGRAM STUDI: D4 TEKNIK INFORMATIKA} \\
\SetCell[c=2]{c} \textbf{RENCANA PEMBELAJARAN SEMESTER (RPS)} & \\
\end{tblr}
\begin{longtblr}{width=\textwidth,colspec={|Q[l,m,1.9cm]|X[0.85,l,m]|X[1.45,l,m]|X[0.75,l,m]|X[1.15,l,m]|},hlines,vlines,hspan=minimal,rowsep=1pt,colsep=2pt,row{1,3}={bg=headgray,font=\bfseries}}
MATA KULIAH & KODE & BOBOT (SKS)/jam & SEMESTER & TGL. PENYUSUNAN \\
Pengembangan Karir & RTI256002 & 2 SKS/4 jam & 6 & 31 Januari 2026 \\
OTORISASI & Dosen Pengembang RPS & Koordinator RMK & \SetCell[c=2]{l} Ka PRODI &  \\
 & Adevian Fairuz Pratama, S.ST, M.Eng &  & \SetCell[c=2]{l}{\vspace{8mm}\par Dr. Ely Setyo Astuti, S.T., M.T.} &  \\
\SetCell[r=13]{l} Capaian Pembelajaran (CP) & \SetCell[c=4]{l,bg=headgray,font=\bfseries} Capaian Pembelajaran Lulusan Program Studi (CPL-Prodi) &  &  &  \\
 & CPL01 & \SetCell[c=3]{l} Mampu menunjukkan nilai ketakwaan, kedisiplinan, profesionalisme, etika profesi, dan komitmen terhadap pembelajaran berkelanjutan dalam konteks sosial dan teknologi. &  &  \\
 & CPL03 & \SetCell[c=3]{l} Mampu mengelola diri, bekerja sama dalam tim, berkomunikasi secara lisan dan tulisan, serta mempresentasikan gagasan secara profesional dalam lingkungan kerja. &  &  \\
 & \SetCell[c=4]{l,bg=headgray,font=\bfseries} Capaian Pembelajaran mata kuliah (CPL-MK) &  &  &  \\
 & CPMK0102 & \SetCell[c=3]{l} Mahasiswa mampu menerapkan nilai ketakwaan, integritas, dan etika profesional sebagai dasar pengambilan keputusan karir dan perilaku profesional di lingkungan kerja. &  &  \\
 & CPMK0303 & \SetCell[c=3]{l} Mahasiswa mampu mengembangkan diri secara berkelanjutan, berkomunikasi secara profesional, dan mengelola karir secara sistematis untuk mencapai tujuan karir jangka panjang. &  &  \\
 & \SetCell[c=4]{l,bg=headgray,font=\bfseries} Kemampuan akhir tiap tahapan belajar (Sub-CPMK) &  &  &  \\
 & SCPMK0102-05301 & \SetCell[c=3]{l} Mahasiswa mampu menerapkan nilai etika profesional dan sikap integritas dalam pengambilan keputusan karir sehari-hari dan lingkungan kerja profesional. &  &  \\
 & SCPMK0303-05301 & \SetCell[c=3]{l} Mahasiswa mampu mengidentifikasi dan menganalisis minat, bakat, potensi diri, serta peluang karir sesuai kompetensi dan tren industri informatika. &  &  \\
 & SCPMK0303-05302 & \SetCell[c=3]{l} Mahasiswa mampu merancang strategi personal branding dan komunikasi profesional yang efektif untuk pengembangan karir di era digital. &  &  \\
 & SCPMK0303-05303 & \SetCell[c=3]{l} Mahasiswa mampu menyusun Career Development Plan (CDP) yang sistematis mencakup tujuan SMART, analisis gap kompetensi, strategi pengembangan diri, dan timeline pencapaian. &  &  \\
\SetCell[r=2]{l} Penilaian & \SetCell[c=4]{l} *Nilai CPMK didapat dari agregasi nilai SUB-CPMK yang dibebankan &  &  &  \\
 & \SetCell[c=4]{l}{\tiny\begin{tblr}{width=\linewidth,colspec={|Q[c,m,0.1\linewidth]|Q[c,m,0.16\linewidth]|X[l,m]|X[l,m]|X[l,m]|X[l,m]|Q[c,m,0.08\linewidth]|},hlines,vlines,hspan=minimal,rowsep=1pt,colsep=0.7pt,row{1,2}={bg=headgray,font=\bfseries}}
Kode CPMK & Kode SCPMK & \SetCell[c=4]{c} Bobot per Bentuk Penilaian & & & & Total Bobot \\
 & & Tugas & UTS & PBL & UAS & \\
\SetCell[r=1]{c} CPMK0102 & SCPMK0102-05301 & 6\%: etika profesi \& integritas karir & 8\%: presentasi portofolio karir & & & 14\% \\
\SetCell[r=3]{c} CPMK0303 & SCPMK0303-05301 & 8\%: minat, potensi \& peluang karir & 9\%: presentasi portofolio karir & & 4\%: evaluasi akhir & 21\% \\
 & SCPMK0303-05302 & 6\%: personal branding \& komunikasi & 8\%: presentasi portofolio karir & & 3\%: evaluasi akhir & 17\% \\
 & SCPMK0303-05303 & 8\%: perencanaan karir \& CDP & & 40\%: Career Development Plan & & 48\% \\
\SetCell[c=6]{r}\textbf{Total Per Penilaian} & & & & & & \textbf{100\%} \\
\end{tblr}} &  &  &  \\
Diskripsi Singkat Mata Kuliah & \SetCell[c=4]{l} Mata kuliah ini bertujuan agar mahasiswa memahami dan mampu mengelola pengembangan karir secara sistematis sehingga dapat mencapai kepuasan dan kesuksesan dalam berkarir. Materi meliputi pemetaan minat dan potensi diri, profesi bidang informatika, personal branding, komunikasi profesional, manajemen dan perencanaan karir, serta penyusunan Career Development Plan (CDP) melalui pendekatan individual project dan presentasi profesional. &  &  &  \\
Bahan Kajian / Materi Pembelajaran & \SetCell[c=4]{l}\begin{enumerate}\item Pengantar pengembangan karir: konsep, tahapan, dan faktor karir\item Pemetaan minat, bakat, dan potensi diri (self-assessment)\item Karir mandiri dan fleksibel di era digital (gig economy, freelance, startup)\item Profesi bidang informatika: klasifikasi, kompetensi, dan jalur karir\item Karir organisasi vs freelance: perbandingan dan strategi\item Karir era Industri 4.0: otomasi, AI, IoT, dan lifelong learning\item Personal branding: identitas digital, LinkedIn, dan reputasi profesional\item Public speaking profesional: komunikasi verbal/nonverbal dan presentasi karir\item Manajemen karir jangka panjang dan networking profesional\item Perencanaan karir: SMART goals, roadmap, dan strategi pencapaian\item Penilaian kinerja: KPI, feedback, dan pengaruh terhadap karir\item Internasionalisasi karir: peluang kerja global dan persyaratan internasional\item Budaya manajerial global dan kolaborasi lintas budaya\item Pembuatan CV profesional yang relevan dan ATS-friendly\end{enumerate} &  &  &  \\
\SetCell[r=2]{l} Pustaka & \SetCell[c=4]{l} \textbf{Utama:}\begin{enumerate}\item Widyanti, R. (2021). Manajemen Karir (Teori, Konsep dan Praktik). Media Sains Indonesia.\end{enumerate} &  &  &  \\
 & \SetCell[c=4]{l} \textbf{Pendukung:}\begin{enumerate}\item Sinambela, L. P. (2021). Manajemen Sumber Daya Manusia: Membangun tim kerja yang solid untuk meningkatkan kinerja. Bumi Aksara.\item Materi ajar/modul dari dosen/tim pengajar dan LMS.\end{enumerate} &  &  &  \\
Media Pembelajaran & \SetCell[c=2]{l} \textbf{Software:} Microsoft Office, LMS, Zoom/platform daring. &  & \SetCell[c=2]{l} \textbf{Hardware:} Gadget/laptop, proyektor, jaringan internet. &  \\
Nama Dosen Pengampu & \SetCell[c=4]{l}\begin{enumerate}\item Agung Nugroho Pramudhita, ST., MT.\item Hendra Pradibta\end{enumerate} &  &  &  \\
Matakuliah Syarat & \SetCell[c=4]{l} -- &  &  &  \\
\end{longtblr}
\begin{longtblr}{width=\textwidth,colspec={|Q[c,m,0.06\textwidth]|X[1.35,l,m]|X[1.08,l,m]|X[1.16,l,m]|X[0.7,l,m]|X[1.24,l,m]|X[1.06,l,m]|X[0.98,l,m]|Q[c,m,0.05\textwidth]|},hlines,vlines,hspan=minimal,rowhead=2,row{1,2}={bg=headgreen,font=\bfseries},rowsep=1pt,colsep=1.5pt}
Mgg.\par Ke & Kemampuan Akhir Yang Direncanakan (Sub-CP-MK) & Bahan kajian (Materi Pembelajaran) & Bentuk dan Metode Pembelajaran & Estimasi Waktu & Pengalaman Belajar Mahasiswa & Kriteria \& Bentuk Penilaian & Indikator Penilaian & Bobot Penilaian (\%) \\
(1) & (2) & (3) & (4) & (5) & (6) & (7) & (8) & (9) \\
1 & SCPMK0102-05301, SCPMK0303-05301 & Pengantar pengembangan karir. & Modalitas: Blended Learning. Bentuk: Luring/Daring. Metode: Ceramah, tanya jawab, diskusi, tugas. & 1 x 2 x 50' tatap muka; 1 x 2 x 50' tugas/praktik mandiri. & Mahasiswa mengerjakan aktivitas terarah untuk pengantar pengembangan karir dan menunjukkan evidence hasil belajar. & Bentuk penilaian: Tugas. Mengacu rubrik RTM. & Ketepatan konsep; kedalaman refleksi; kebaruan ide dan gagasan. & 2 \\
2 & SCPMK0303-05301 & Pemetaan minat dan potensi karir. & Modalitas: Blended Learning. Bentuk: Luring/Daring. Metode: Ceramah, studi literatur, diskusi, tugas. & 1 x 2 x 50' tatap muka; 1 x 2 x 50' tugas/praktik mandiri. & Mahasiswa mengerjakan aktivitas terarah untuk pemetaan minat dan potensi karir serta menunjukkan evidence hasil belajar. & Bentuk penilaian: Tugas. Mengacu rubrik RTM. & Ketepatan konsep; kedalaman refleksi; kebaruan ide dan gagasan. & 2 \\
3 & SCPMK0303-05301 & Karir mandiri dan fleksibel di era digital. & Modalitas: Blended Learning. Bentuk: Luring/Daring. Metode: Ceramah, tanya jawab, diskusi, tugas. & 1 x 2 x 50' tatap muka; 1 x 2 x 50' tugas/praktik mandiri. & Mahasiswa mengerjakan aktivitas terarah untuk karir mandiri di era digital serta menunjukkan evidence hasil belajar. & Bentuk penilaian: Tugas. Mengacu rubrik RTM. & Ketepatan konsep; kedalaman refleksi; kebaruan ide dan gagasan. & 2 \\
4 & SCPMK0303-05301 & Profesi bidang informatika. & Modalitas: Blended Learning. Bentuk: Luring/Daring. Metode: Ceramah, studi literatur, diskusi, tugas. & 1 x 2 x 50' tatap muka; 1 x 2 x 50' tugas/praktik mandiri. & Mahasiswa mengerjakan aktivitas terarah untuk profesi bidang informatika serta menunjukkan evidence hasil belajar. & Bentuk penilaian: Tugas. Mengacu rubrik RTM. & Ketepatan konsep; kedalaman refleksi; kebaruan ide dan gagasan. & 2 \\
5 & SCPMK0303-05301 & Karir organisasi vs freelance. & Modalitas: Blended Learning. Bentuk: Luring/Daring. Metode: Ceramah, tanya jawab, diskusi, tugas. & 1 x 2 x 50' tatap muka; 1 x 2 x 50' tugas/praktik mandiri. & Mahasiswa mengerjakan aktivitas terarah untuk perbandingan karir organisasi dan freelance serta menunjukkan evidence hasil belajar. & Bentuk penilaian: Tugas. Mengacu rubrik RTM. & Ketepatan konsep; kedalaman refleksi; kebaruan ide dan gagasan. & 2 \\
6 & SCPMK0303-05301 & Karir era Industri 4.0. & Modalitas: Blended Learning. Bentuk: Luring/Daring. Metode: Ceramah, diskusi, tugas. & 1 x 2 x 50' tatap muka; 1 x 2 x 50' tugas/praktik mandiri. & Mahasiswa mengerjakan aktivitas terarah untuk karir di era Industri 4.0 serta menunjukkan evidence hasil belajar. & Bentuk penilaian: Tugas. Mengacu rubrik RTM. & Ketepatan konsep; kedalaman refleksi; kebaruan ide dan gagasan. & 2 \\
7 & SCPMK0303-05302 & Personal branding profesional. & Modalitas: Blended Learning. Bentuk: Luring/Daring. Metode: Ceramah, diskusi, tugas. & 1 x 2 x 50' tatap muka; 1 x 2 x 50' tugas/praktik mandiri. & Mahasiswa mengerjakan aktivitas terarah untuk personal branding profesional serta menunjukkan evidence hasil belajar. & Bentuk penilaian: Tugas. Mengacu rubrik RTM. & Ketepatan konsep; kedalaman refleksi; kebaruan ide dan gagasan. & 2 \\
8 & SCPMK0102-05301, SCPMK0303-05301, SCPMK0303-05302 & UTS: presentasi portofolio karir awal. & Modalitas: Blended Learning. Bentuk: Luring/Daring. Metode: Presentasi individual. & 1 x 2 x 50' tatap muka; 1 x 2 x 50' tugas/praktik mandiri. & Mahasiswa mempresentasikan portofolio karir awal yang mencakup analisis diri, profil karir, dan strategi personal branding. & Bentuk penilaian: UTS presentasi portofolio karir. Mengacu rubrik RTM. & Kelengkapan portofolio; kemampuan presentasi; argumentasi rencana karir. & 25 \\
9 & SCPMK0303-05302 & Public speaking profesional. & Modalitas: Blended Learning. Bentuk: Luring/Daring. Metode: Ceramah, diskusi, tugas. & 1 x 2 x 50' tatap muka; 1 x 2 x 50' tugas/praktik mandiri. & Mahasiswa mengerjakan aktivitas terarah untuk public speaking profesional serta menunjukkan evidence hasil belajar. & Bentuk penilaian: Tugas. Mengacu rubrik RTM. & Ketepatan konsep; kedalaman refleksi; kebaruan ide dan gagasan. & 2 \\
10 & SCPMK0303-05303 & Manajemen karir jangka panjang. & Modalitas: Blended Learning. Bentuk: Luring/Daring. Metode: Ceramah, diskusi, tugas. & 1 x 2 x 50' tatap muka; 1 x 2 x 50' tugas/praktik mandiri. & Mahasiswa mengerjakan aktivitas terarah untuk manajemen karir jangka panjang serta menunjukkan evidence hasil belajar. & Bentuk penilaian: Tugas. Mengacu rubrik RTM. & Ketepatan konsep; kedalaman refleksi; kebaruan ide dan gagasan. & 2 \\
11 & SCPMK0303-05303 & Perencanaan karir: SMART goals dan roadmap. & Modalitas: Blended Learning. Bentuk: Luring/Daring. Metode: Ceramah, diskusi, tugas. & 1 x 2 x 50' tatap muka; 1 x 2 x 50' tugas/praktik mandiri. & Mahasiswa mengerjakan aktivitas terarah untuk perencanaan karir dengan SMART goals dan roadmap serta menunjukkan evidence hasil belajar. & Bentuk penilaian: Tugas. Mengacu rubrik RTM. & Ketepatan konsep; kedalaman refleksi; kebaruan ide dan gagasan. & 2 \\
12 & SCPMK0303-05303 & Penilaian kinerja dan KPI. & Modalitas: Blended Learning. Bentuk: Luring/Daring. Metode: Ceramah, diskusi, tugas. & 1 x 2 x 50' tatap muka; 1 x 2 x 50' tugas/praktik mandiri. & Mahasiswa mengerjakan aktivitas terarah untuk penilaian kinerja dan KPI serta menunjukkan evidence hasil belajar. & Bentuk penilaian: Tugas. Mengacu rubrik RTM. & Ketepatan konsep; kedalaman refleksi; kebaruan ide dan gagasan. & 2 \\
13 & SCPMK0303-05303 & Internasionalisasi karir. & Modalitas: Blended Learning. Bentuk: Luring/Daring. Metode: Ceramah, diskusi, tugas. & 1 x 2 x 50' tatap muka; 1 x 2 x 50' tugas/praktik mandiri. & Mahasiswa mengerjakan aktivitas terarah untuk internasionalisasi karir serta menunjukkan evidence hasil belajar. & Bentuk penilaian: Tugas. Mengacu rubrik RTM. & Ketepatan konsep; kedalaman refleksi; kebaruan ide dan gagasan. & 2 \\
14 & SCPMK0102-05301 & Budaya manajerial global dan etika lintas budaya. & Modalitas: Blended Learning. Bentuk: Luring/Daring. Metode: Ceramah, diskusi, tugas. & 1 x 2 x 50' tatap muka; 1 x 2 x 50' tugas/praktik mandiri. & Mahasiswa mengerjakan aktivitas terarah untuk budaya manajerial global dan etika lintas budaya serta menunjukkan evidence hasil belajar. & Bentuk penilaian: Tugas. Mengacu rubrik RTM. & Ketepatan konsep; kedalaman refleksi; kebaruan ide dan gagasan. & 2 \\
15 & SCPMK0303-05302 & Pembuatan CV profesional dan ATS-friendly. & Modalitas: Blended Learning. Bentuk: Luring/Daring. Metode: Ceramah, diskusi, tugas. & 1 x 2 x 50' tatap muka; 1 x 2 x 50' tugas/praktik mandiri. & Mahasiswa mengerjakan aktivitas terarah untuk pembuatan CV profesional serta menunjukkan evidence hasil belajar. & Bentuk penilaian: Tugas. Mengacu rubrik RTM. & Ketepatan konsep; kedalaman refleksi; kebaruan ide dan gagasan. & 2 \\
16 & SCPMK0102-05301--SCPMK0303-05303 & PBL: presentasi Career Development Plan. & Modalitas: Blended Learning. Bentuk: Luring/Daring. Metode: Presentasi proyek individual, peer review. & 1 x 2 x 50' tatap muka; 1 x 2 x 50' tugas/praktik mandiri. & Mahasiswa mempresentasikan Career Development Plan (CDP) secara menyeluruh yang mencakup tujuan karir, analisis gap kompetensi, strategi pengembangan, dan timeline. & Bentuk penilaian: PBL Career Development Plan. Mengacu rubrik RTM. & Kelengkapan CDP; kedalaman analisis; kualitas presentasi dan pertahanan argumen. & 40 \\
17 & SCPMK0303-05301, SCPMK0303-05302 & UAS: evaluasi akhir pengembangan karir. & Modalitas: Blended Learning. Bentuk: Luring/Daring. Metode: Evaluasi reflektif dan presentasi singkat. & 1 x 2 x 50' tatap muka; 1 x 2 x 50' tugas/praktik mandiri. & Mahasiswa melakukan evaluasi akhir dan merefleksikan seluruh proses pengembangan karir selama semester, serta menunjukkan konsistensi nilai etika profesional. & Bentuk penilaian: UAS evaluasi akhir pengembangan karir. Mengacu rubrik RTM. & Kualitas refleksi akhir; konsistensi etika profesional. & 7 \\
\end{longtblr}
